\chapter{Identity, Passwords, and Access Risk in the Age of Agentic AI}

\begin{learningobjectives}
By the end of this chapter, the reader should be able to:
\begin{itemize}
    \item distinguish authentication, authorization, and accounting in financial institutions;
    \item explain why identity is one of the most important control layers in modern financial systems;
    \item understand password, session, privilege, and service-account risk without relying on offensive techniques;
    \item evaluate multifactor authentication, least privilege, segregation of duties, and privileged-access controls;
    \item describe, at a safe and non-operational level, the AI models, agentic architectures, and automation strategies that may increase identity-related cyber threats;
    \item explain why AI agents themselves must be governed as institutional identities with limited authority;
    \item connect identity compromise to payments, trading, asset management, treasury, custody, and financial-market infrastructure.
\end{itemize}
\end{learningobjectives}

\section{Identity Is a Financial Control}

A financial institution is full of decisions that depend on identity.

Who may enter an order? Who may change a risk limit? Who may release a payment? Who may access client data? Who may modify a pricing source? Who may deploy software? Who may approve a model change? Who may instruct a custodian?

In each case, the institution must answer three questions:

\begin{equation}
\text{Who are you?}
\qquad
\text{What may you do?}
\qquad
\text{What evidence remains?}
\end{equation}

These correspond to three classical control concepts:

\begin{itemize}
    \item \textbf{Authentication}: establishing the identity of the user or system.
    \item \textbf{Authorization}: determining which actions that identity is permitted to perform.
    \item \textbf{Accounting or auditability}: recording what the identity actually did.
\end{itemize}

For finance students, this should feel familiar. Financial businesses already depend on authority structures. A junior trader and the head of trading do not have the same limits. A payment preparer and payment approver should not have identical authority. An analyst who prepares a valuation should not necessarily have the authority to approve the final NAV. A portfolio manager may propose a transaction, while operations personnel confirm and settle it.

Cybersecurity extends these control principles into digital systems.

The central risk is not simply that an unauthorized person obtains a password. The real risk is that the institution mistakenly grants authority to someone or something that should not possess it.

\begin{intuition}
Identity risk is the risk that a financial institution assigns trust, authority, or access to the wrong human, device, service, or autonomous agent.
\end{intuition}

\section{Authentication Is Not Authorization}

Authentication and authorization are often confused.

Suppose Alice successfully logs into a portfolio-management system. Authentication answers:

\begin{quote}
Is this really Alice?
\end{quote}

Authorization answers:

\begin{quote}
Even if this is Alice, is Alice permitted to change the portfolio's benchmark, export the complete client list, or modify settlement instructions?
\end{quote}

A secure financial system needs both.

Strong authentication combined with excessive authorization remains dangerous. If a legitimate employee has far more permissions than required, the institution has created unnecessary operational risk. The account may be misused intentionally, compromised externally, or used incorrectly by accident.

Conversely, perfect authorization rules are useless if authentication is weak. The system may have carefully designed roles, but if an unauthorized actor can impersonate a highly privileged user, those roles no longer protect the institution.

A useful conceptual decomposition is:

\begin{equation}
R_{\text{identity}}
=
P(\text{identity misuse})
\times
I(\text{authority attached to identity}).
\end{equation}

The second term is frequently underestimated.

A compromised account that can read a public research archive is different from a compromised account that can release treasury payments.

\section{Why Passwords Became a Weak Control}

Passwords are historically convenient because they allow a person to prove knowledge of a secret. But passwords have several structural weaknesses.

Users may choose weak secrets. They may reuse the same password across services. They may reveal it accidentally. A password may be exposed in a breach at another organization. Users may respond to convincing fraudulent messages. Shared passwords may become impossible to attribute to one person.

Finance professionals do not need to know how password attacks are technically executed. They need to understand why a password alone may provide insufficient assurance for high-impact actions.

The weakness becomes particularly important when the same identity has broad economic authority.

Consider a treasury employee whose account can prepare a payment. If the account is protected only by a password, the institution is making a strong assumption: knowledge of one string is sufficient evidence that a high-value financial action comes from the correct human.

That is a weak governance proposition.

\section{Multifactor Authentication}

Multifactor authentication improves assurance by requiring more than one type of evidence.

The usual categories are:

\begin{itemize}
    \item something the user knows;
    \item something the user possesses;
    \item something the user is.
\end{itemize}

The objective is not to make authentication inconvenient. It is to make one failure insufficient.

For a financial institution, multifactor authentication should be especially important for:

\begin{itemize}
    \item privileged administrator accounts;
    \item payment and treasury systems;
    \item remote access;
    \item cloud administration;
    \item sensitive client data;
    \item production software deployment;
    \item trading administration;
    \item critical vendor interfaces.
\end{itemize}

However, multifactor authentication is not a complete solution. An authenticated session can still be misused. A user can authorize the wrong action. A malicious process may operate inside an already authenticated environment. A poorly designed recovery procedure may bypass the stronger control.

This is why identity defense requires a chain of controls rather than one mechanism.

\section{Least Privilege}

The principle of least privilege says that an identity should possess only the authority required to perform its legitimate function.

This principle has a direct analogue in financial control.

A trader may need authority to enter orders within defined limits but not to change the limits themselves.

An operations employee may need to prepare a payment but not approve it.

A valuation analyst may need access to pricing data but not the ability to alter production source files.

An AI research assistant may need read access to research documents but not permission to submit trades.

Least privilege reduces the impact of identity compromise.

If an account is compromised but its permissions are narrow, the potential damage remains constrained.

\begin{equation}
\text{Compromise}
\times
\text{Limited Authority}
\rightarrow
\text{Limited Loss Potential}.
\end{equation}

The reverse is equally important:

\begin{equation}
\text{Compromise}
\times
\text{Broad Authority}
\rightarrow
\text{Potentially Large Loss}.
\end{equation}

\section{Segregation of Duties}

Segregation of duties prevents one identity from controlling an entire high-risk process.

A classic financial example is payment processing:

\begin{equation}
\text{Prepare}
\neq
\text{Approve}
\neq
\text{Release}.
\end{equation}

The same principle applies to cybersecurity and AI.

A developer who writes production code should not necessarily have unilateral authority to deploy it.

An employee who creates a new beneficiary should not necessarily be allowed to release a large transfer to that beneficiary.

An AI agent that recommends a trade should not automatically possess the authority to execute it.

Segregation of duties assumes that individual controls can fail. The architecture is designed so that one failure does not automatically create the final adverse outcome.

\section{Privileged Identities}

Some accounts deserve special treatment because they can change the environment itself.

Privileged identities may be able to:

\begin{itemize}
    \item create or delete users;
    \item change permissions;
    \item access confidential databases;
    \item modify system configuration;
    \item disable monitoring;
    \item deploy software;
    \item access backups;
    \item alter security controls.
\end{itemize}

In a financial institution, privileged-access management should therefore be viewed as a concentration-risk problem.

If one identity can modify many critical systems, the institution has concentrated operational power.

The control response may include stronger authentication, temporary privilege elevation, independent approval, continuous monitoring, session recording, and rapid revocation.

\section{Machine Identities and Service Accounts}

Not every identity represents a human.

Financial institutions use service accounts, application identities, API credentials, software certificates, automated jobs, and machine-to-machine permissions.

These identities can be extremely powerful because they operate continuously and may have broad access.

Historically, organizations often gave less attention to machine identities than human identities. That becomes increasingly dangerous as automation expands.

The arrival of autonomous AI agents makes the issue much more important.

An agent that can read email, query databases, call APIs, edit documents, execute code, or initiate workflows is effectively an institutional user.

It should therefore have an identity.

It should have a role.

It should have permissions.

It should have limits.

It should leave an audit trail.

\section{The AI Threat Lens: Models Are Becoming Operational Actors}

The identity problem changes materially when frontier AI models are connected to tools.

OpenAI's current agent architecture describes agents as systems built from models, tools, state or memory, and orchestration. Its agent tooling also supports handoffs among agents, guardrails, and tracing. Anthropic similarly documents tool-using and managed agent systems in which a model can autonomously invoke approved functions.

This matters because the model is no longer limited to generating text.

A tool-enabled model can potentially:

\begin{equation}
\text{Read}
\rightarrow
\text{Reason}
\rightarrow
\text{Select Tool}
\rightarrow
\text{Act}
\rightarrow
\text{Observe Result}
\rightarrow
\text{Continue}.
\end{equation}

The cyber-risk significance depends on what tools are available.

A model with access only to a public research database has limited authority.

A model with access to internal email, file systems, customer information, payment APIs, code repositories, cloud consoles, and administrative systems has much greater operational power.

The relevant unit of risk is therefore not the model alone.

\begin{equation}
R_{\text{AI system}}
=
f(
\text{model capability},
\text{tools},
\text{permissions},
\text{memory},
\text{autonomy},
\text{coordination}
).
\end{equation}

\section{Frontier Cyber Models}

The frontier is moving rapidly.

OpenAI reports that GPT-5.6 is its strongest cybersecurity model to date and that GPT-5.6 Sol extends capability on long-horizon security tasks. OpenAI's deployment safety analysis also states that these models represent a meaningful increase in cybersecurity capability, while not reaching its highest classified cyber-risk level.

Anthropic has taken a similarly cautious approach with the Mythos family. It describes Claude Mythos as a highly capable cybersecurity model and has limited access to vetted partners. Anthropic's research around Mythos Preview reports substantial gains in computer-security tasks and explicitly frames those gains as useful for defensive vulnerability discovery while acknowledging misuse risk.

These examples illustrate a general trend rather than a competition among model brands:

\begin{intuition}
The amount of expert-level cyber reasoning that can be automated is increasing, and the time required to perform complex analysis is declining.
\end{intuition}

For financial institutions, this means that historical assumptions about the scarcity of sophisticated attackers may become less reliable.

\section{Typical AI Models Used in an Attack Architecture}

This book will not describe operational attack procedures. It is nevertheless important for risk managers to understand the types of AI components that can appear in a malicious architecture.

A high-level architecture may include several model roles.

\subsection{Reasoning Model}

A frontier reasoning model can interpret evidence, maintain a plan, compare alternatives, and revise its approach.

Its risk contribution is persistence and analytical scale.

\subsection{Coding Model}

A coding-specialized model can analyze software, generate or modify code, and help automate technical tasks.

Its risk contribution is lowering the cost of software analysis and experimentation.

\subsection{Vision or Multimodal Model}

A multimodal model can interpret screenshots, dashboards, interfaces, diagrams, and documents.

Its risk contribution is expanding automation beyond purely text-based environments.

\subsection{Small Local Model}

A smaller model may perform repetitive classification, extraction, or coordination tasks cheaply and at scale.

The lesson is important: a threat architecture does not necessarily require the most capable model for every task.

\subsection{Coordinator Model}

A coordinating agent can allocate work among specialized agents, maintain shared state, and decide which task should occur next.

Its risk contribution is orchestration.

\section{Typical Agentic Architectures Used by Threat Actors}

At a safe level of abstraction, several architectures are particularly relevant.

\subsection{Single-Agent Tool Loop}

The simplest architecture is:

\begin{equation}
\text{Observe}
\rightarrow
\text{Plan}
\rightarrow
\text{Tool}
\rightarrow
\text{Evaluate}
\rightarrow
\text{Repeat}.
\end{equation}

The model maintains one objective and iteratively uses available tools.

The defensive implication is that malicious activity may become more persistent and less dependent on continuous human attention.

\subsection{Planner--Executor Architecture}

One agent creates a plan while another performs constrained tasks.

\begin{equation}
\text{Planner}
\rightarrow
\text{Executor}
\rightarrow
\text{Verifier}.
\end{equation}

The verifier checks whether the objective was achieved and sends the task back for revision if necessary.

This structure can increase reliability.

\subsection{Specialist Agent Team}

A coordinated team may include:

\begin{itemize}
    \item an information-gathering agent;
    \item a software-analysis agent;
    \item an identity-analysis agent;
    \item a coordinator;
    \item an evaluator.
\end{itemize}

The risk is not that each specialist is extraordinary. The risk is that parallel agents can divide work and maintain persistent attention.

\subsection{Hierarchical Agent Architecture}

A supervisory agent may manage many lower-level agents.

\begin{equation}
\text{Supervisor}
\rightarrow
\left\{
A_1,A_2,\ldots,A_n
\right\}
\rightarrow
\text{Aggregated Result}.
\end{equation}

This creates scale.

\subsection{Memory-Augmented Agent}

An agent can preserve findings, previous attempts, system information, and workflow lessons across many iterations.

Memory reduces repeated work.

For defenders, this means that an automated adversary may accumulate context over time rather than beginning every interaction from zero.

\section{Typical AI-Enabled Threat Strategies}

Again, the purpose here is threat recognition rather than operational instruction.

The important strategies include:

\subsection{Automation of Repetitive Investigation}

AI can reduce the human effort required to analyze large numbers of systems, configurations, messages, or code artifacts.

The strategic effect is scale.

\subsection{Adaptive Decision-Making}

Instead of following a rigid script, an agentic system can evaluate outcomes and change its next action.

The strategic effect is adaptability.

\subsection{Parallelization}

Multiple agents can work simultaneously.

The strategic effect is lower marginal cost per target or task.

\subsection{Persistent Operation}

Agents can continue working for extended periods without normal human attention limits.

The strategic effect is endurance.

\subsection{Personalization}

AI can analyze available context and produce more tailored messages, requests, or interactions.

For financial institutions, this increases the importance of independent verification procedures.

\subsection{Tool Chaining}

The greatest risk often comes from combining several legitimate capabilities.

A model may read information with one tool, transform it with another, and initiate a workflow with a third.

The individual tools may all be legitimate. The risk emerges from the sequence.

\section{Identity Attacks in an Agentic World}

AI changes the identity threat in three ways.

First, attackers may automate the analysis of organizational identities and relationships.

Second, AI can make fraudulent communications more context-aware, linguistically convincing, and persistent.

Third, compromised legitimate identities may themselves be used by autonomous systems.

This makes identity assurance more important, not less.

A financial institution should increasingly assume that a convincing email, voice message, document, or chat interaction is not sufficient proof of identity.

Independent verification matters.

\section{The New Control: Agent Identity}

Financial institutions increasingly need an explicit concept of \textbf{agent identity}.

An autonomous agent should not inherit the full authority of the human who created it.

Instead:

\begin{equation}
\text{Human Identity}
\neq
\text{Agent Identity}.
\end{equation}

The institution should be able to answer:

\begin{itemize}
    \item Which agent performed the action?
    \item Which human or business function owns the agent?
    \item Which model powers it?
    \item Which tools can it call?
    \item Which data can it read?
    \item Which actions can it execute?
    \item What monetary or operational limits apply?
    \item When does its authorization expire?
    \item Can its authority be revoked immediately?
\end{itemize}

This is the agentic equivalent of access governance.

\section{A Finance Example: Autonomous Treasury}

Suppose a financial institution deploys an autonomous treasury agent.

The initial version can read cash balances, forecast liquidity, and propose transfers.

This is a decision-support system.

Management then adds the ability to submit transfer instructions.

The agent has become an operational actor.

Later, management considers allowing automatic execution below a threshold.

At that point, the agent possesses economic authority.

A sound control structure could separate:

\begin{equation}
\text{Forecast Agent}
\rightarrow
\text{Recommendation Agent}
\rightarrow
\text{Payment Preparation}
\rightarrow
\text{Independent Approval}.
\end{equation}

The institution should be cautious about collapsing these functions into one autonomous identity.

\begin{risklens}
The central financial principle is the same whether the actor is human or artificial: authority should be proportional to need, economically bounded, independently observable, and revocable.
\end{risklens}

\section{A Trading Example: The Autonomous Research Agent}

Consider an investment manager using an AI research agent.

The agent can read filings, internal research, portfolio information, and market data.

That architecture already raises confidentiality questions.

Suppose the agent is later permitted to generate trade recommendations.

The risk increases, but the agent still has no execution authority.

Suppose it is finally connected directly to an order-management system.

The architecture has crossed a major governance boundary.

The institution must now consider:

\begin{itemize}
    \item position limits;
    \item instrument restrictions;
    \item pre-trade risk checks;
    \item human approval thresholds;
    \item model drift;
    \item prompt or context manipulation;
    \item kill switches;
    \item independent trade surveillance;
    \item complete action logs.
\end{itemize}

The relevant principle is not ``never use autonomous agents.'' It is ``never confuse analytical capability with authorized economic authority.''

\section{Defensive Strategies for the AI Identity Era}

The defensive response should combine conventional identity controls with agent-specific governance.

A pragmatic architecture includes:

\begin{enumerate}
    \item strong human authentication;
    \item unique identities for agents and services;
    \item least privilege;
    \item short-lived credentials where possible;
    \item segregation of duties;
    \item transaction and action limits;
    \item independent logging;
    \item continuous behavioral monitoring;
    \item explicit approval for high-impact actions;
    \item immediate revocation capability.
\end{enumerate}

Financial institutions should also avoid embedding long-lived credentials directly into AI prompts, code, notebooks, or agent memory.

The distinction between reasoning state and secret material should be explicit.

\section{Mini Case: The Autonomous Payment Assistant}

A fictional bank deploys an AI payment assistant.

Phase 1 allows the agent to read payment requests and classify them.

Phase 2 allows it to prepare payments for human approval.

Phase 3 allows it to release payments below \$50,000.

At Phase 1, the principal risks concern confidentiality and classification error.

At Phase 2, authorization risk increases because the agent can create actionable financial instructions.

At Phase 3, the agent can directly create financial loss.

The bank should therefore reassess the control architecture at every phase.

The appropriate questions include:

\begin{itemize}
    \item Is the agent represented by a unique machine identity?
    \item Are its privileges different from those of its human owner?
    \item Can it create new beneficiaries?
    \item Is the \$50,000 limit per transaction or aggregate?
    \item Can many small transfers bypass the economic intent of the limit?
    \item Is a second control independent of the agent?
    \item Are all actions recorded outside the agent's own environment?
    \item Can operations disable the agent immediately?
\end{itemize}

This is a practical example of cybersecurity, operational risk, model risk, and financial control converging.

\section{Safe Notebook Lab}

\begin{notebooklab}
The companion notebook for this chapter should remain entirely synthetic.

Students can create fictional identities representing:

\begin{itemize}
    \item a trader;
    \item a treasury employee;
    \item an administrator;
    \item a research agent;
    \item a payment agent.
\end{itemize}

Each identity receives a transparent permission set.

The notebook can then simulate harmless requests such as ``read portfolio,'' ``prepare payment,'' or ``change limit'' and determine whether each identity is authorized.

A second exercise can compare three control architectures:
\begin{enumerate}
    \item one broad privileged identity;
    \item segregated human roles;
    \item segregated human and AI-agent identities.
\end{enumerate}

The educational objective is to quantify how authority concentration affects potential loss.

The notebook should not use real credentials, external systems, authentication services, payment APIs, trading APIs, or network scanning.
\end{notebooklab}

\section{A Pragmatic Identity Checklist for Finance Professionals}

When evaluating an important financial process, ask:

\begin{enumerate}
    \item Who can perform the action?
    \item How is each identity authenticated?
    \item What is the maximum economic authority of each identity?
    \item Is approval independent from initiation?
    \item Which machine identities and AI agents participate?
    \item Are agent permissions narrower than the permissions of their human owners?
    \item Can one identity change its own permissions?
    \item Are high-impact actions independently logged?
    \item How quickly can access be revoked?
    \item What happens if the identity is correct but the decision is wrong?
\end{enumerate}

The final question is especially important in the AI era. Authentication can prove that the correct agent performed the action. It does not prove that the action was economically sensible.

\section{Chapter Summary}

The central lessons are:

\begin{enumerate}
    \item Identity is a financial control because digital identities exercise economic authority.
    \item Authentication establishes identity; authorization defines permitted actions; accounting preserves evidence.
    \item Passwords alone provide weak assurance for high-impact financial processes.
    \item Multifactor authentication, least privilege, segregation of duties, and privileged-access controls reduce the probability and impact of identity misuse.
    \item Machine identities are as important as human identities.
    \item Autonomous AI agents should be treated as institutional users with unique identities, explicit permissions, limits, ownership, and audit trails.
    \item Frontier models such as GPT-5.6-class systems and Anthropic's Mythos family illustrate the rapid increase in automated cybersecurity capability.
    \item Typical malicious AI architectures may involve reasoning models, coding models, coordinators, memory, tool loops, planner--executor patterns, and specialist agent teams.
    \item AI increases threat scale through automation, adaptability, parallelization, persistence, personalization, and tool chaining.
    \item The core control principle remains simple: capability should never imply authority.
\end{enumerate}

\section{Review Questions}

\begin{enumerate}
    \item What is the difference between authentication and authorization?
    \item Why is identity risk fundamentally a financial-control issue?
    \item Why does least privilege reduce cyber loss severity?
    \item How does segregation of duties protect payment and trading processes?
    \item Why should service accounts and machine identities receive the same governance attention as human identities?
    \item What additional risks arise when a language model is connected to tools?
    \item Describe the difference between a single-agent tool loop and a planner--executor architecture.
    \item Why can a multi-agent team create greater operational risk than one model acting alone?
    \item What does the statement ``capability should never imply authority'' mean for an autonomous trading or payment agent?
    \item In the payment-assistant case, which controls would you require before allowing any autonomous payment execution?
\end{enumerate}
